\section{The source's conjugate
families}\label{the-sources-conjugate-families}

First dispatch the small index sets in source Theorem 7.3. For a single
freely independent generator, the group is infinite cyclic and its sole
omitted-generator subgroup is trivial. Its singleton cosets form a
one-type geometry: every chamber is one object and translation acts
regularly. There is no residue of rank at least two to check. If the
empty index set is allowed, the empty system has the single empty
chamber and trivial group. Thus these endpoints satisfy the source
flag-transitivity assertion; the following nontrivial family uses the
common theory for rank at least two.

Let \(F_n=F(x_1,\ldots,x_n)\) with \(n\ge2\), and put \[
b_{j,\epsilon,i}=x_j^\epsilon x_i x_j^{-\epsilon},
\qquad j\ne i,\quad\epsilon\in\{1,-1\}.
\tag{F9}
\] There are \(r=2n(n-1)\) displayed words. Let \(G\) be the subgroup
they generate.

\textbf{Lemma R2.Core.} The words (F9) freely generate \(G\).

\textbf{Proof.} Construct a finite labeled graph with one central vertex
\(o\) and \(2n\) outer vertices \(v_{j,\epsilon}\). Join \(o\) to
\(v_{j,\epsilon}\) by an oriented edge labeled \(x_j^\epsilon\). At
\(v_{j,\epsilon}\) put a loop labeled \(x_i\) for every \(i\ne j\). Map
this graph to the rose labeled \(x_1,\ldots,x_n\).

At each vertex, all outgoing oriented labels are different: at \(o\)
they are the \(2n\) signed letters; at \(v_{j,\epsilon}\) the edge back
has label \(x_j^{-\epsilon}\), and the loop directions have labels
\(x_i^{\pm1}\) for \(i\ne j\). Thus a reduced edge path maps to a
reduced word. A nonempty reduced closed path cannot map to the identity
in \(F_n\), so the induced homomorphism on fundamental groups is
injective. The spokes form a spanning tree. The associated free basis
consists of a spoke, one outer loop, and the reverse spoke, giving
exactly (F9). Equivalently the graph rank is \[
\bigl(2n+2n(n-1)\bigr)-(1+2n)+1=2n(n-1).
\] This proves the assertion. \(\square\)

The finite graph is the core relevant to the displayed finite generating
set. Attaching loops at every integer vertex of entire coordinate axes,
as a literal reading of the subsequent source description suggests,
would instead give infinitely many independent loops. No such enlarged
graph is used here. Trees may be attached to the finite core to obtain
the subgroup covering without altering its fundamental group.

\textbf{Theorem R2.SourceFamily.} Form the source coset system on \(G\)
with one maximal standard parabolic omitting each word (F9). Then it is
flag-transitive and residually connected, \(G\) acts faithfully and
regularly on chambers, and its full type quotient is \[
\operatorname{Sym}(2n(n-1)).
\tag{F10}
\] The signed permutation group \(C_2^n\rtimes S_n\) of the ambient
basis acts faithfully by correlations. Its action is a specified
subgroup, not the complete type image.

\textbf{Proof.} Lemma R2.Core identifies \(G\) with the free group on
the displayed basis. Theorem R2.Geometry and Theorem R2.TypeImage apply.
An ambient signed permutation \(x_k\mapsto x_{p(k)}^{\delta_k}\) sends
\[
b_{j,\epsilon,i}\longmapsto
b_{p(j),\epsilon\delta_j,p(i)}^{\delta_i}.
\tag{F11}
\] It therefore permutes the omitted-generator parabolics and acts by
(F6). Its palette action is faithful: fixing all indices
\((j,\epsilon,i)\) forces \(p\) to fix every symbol and forces every
\(\delta_j=1\), since for each \(j\) another index \(i\) exists.
\(\square\)

For \(n=2\), the four words are precisely those printed before Theorem
7.5. The specified ambient group acts on their four types as the square
group of order eight. But an arbitrary three-cycle of the four free
generators also acts by a correlation, and such a permutation is not in
that square group. The full quotient in (F10) is \(S_4\) of order
twenty-four. Thus the quotient assertion in Theorem 7.5 is false; it is
not repaired merely by giving a more detailed proof of the action of the
chosen subgroup.

The geometric conclusions and the chosen finite actions of Theorems
7.5--7.6 are recovered. Their weak-representation conclusions also hold:
the displayed subgroup is a finite-rank free group, so Corollary
UR.WeakFree embeds its full automorphism group into the type-preserving
group, compatibly with its inner subgroup. This supplies a different
inner action from the source's chamber-regular construction. The source's
argument from the chosen finite signed action does not itself establish
the weak conclusion, and the natural inner embedding cannot extend by
Corollary R2.NaturalObstruction. Both distinctions are essential; the
literal existence conclusion is not refuted by the natural obstruction.
